% Auto-gerado por C11_build_T5_latex.py
\begin{table}[htbp]
\centering
\caption{Efeito do timing da observa\c{c}\~ao de $t+1$ sobre as taxas de fluxo \emph{within-person}, Brasil, 2012--2023.}
\label{tab:t5_efeito_ferias}
\begin{threeparttable}
\small
\begin{tabular}{lcccccccc}
\toprule
& & \multicolumn{3}{c}{Medido em Q1 de $t+1$} & \multicolumn{3}{c}{Medido em Q2--Q4 de $t+1$} & $\Delta$ Rep. \\
\cmidrule(lr){3-5} \cmidrule(lr){6-8}
Macroetapa & N & Prom. & Rep. & Evas. & Prom. & Rep. & Evas. & (pp) \\
\midrule
EF iniciais (1\textsuperscript{o}--5\textsuperscript{o} ano) & 195.064 &  62.1 &  29.5 &   0.0 &  64.8 &  25.1 &   0.0 &  -4.3 \\
EF finais (6\textsuperscript{o}--9\textsuperscript{o} ano) & 155.981 &  58.6 &  30.4 &   0.0 &  61.0 &  26.6 &   0.0 &  -3.9 \\
Ensino M\'edio (1\textsuperscript{o}--3\textsuperscript{o} ano) & 73.762 &  56.5 &  37.7 &   0.0 &  60.0 &  33.5 &   0.0 &  -4.3 \\
\bottomrule
\end{tabular}
\begin{tablenotes}\footnotesize
\item \textit{Fonte:} PNADC trimestral, painel longitudinal, 2012--2023. Subamostra \emph{within-person}: alunos com observa\c{c}\~ao tanto no Q1 quanto em algum Q$\geq$2 de $t+1$ (cerca de 65\% dos casos com obs.\ em Q1 de $t+1$).
\item \textit{Notas:} A medi\c{c}\~ao em Q1 (jan--mar) coincide com o per\'iodo de f\'erias e in\'icio do ano letivo, em que parte das fam\'ilias ainda reporta a s\'erie anterior. Quando o mesmo aluno \'e remedido em Q2--Q4 (abr--dez, j\'a com a s\'erie nova consolidada), a repet\^encia cai $\sim$5pp e a promo\c{c}\~ao sobe $\sim$3--4pp. Evas\~ao $= 0$ nesta subamostra por constru\c{c}\~ao (exige matr\'icula em ambos os momentos).
\item Coluna $\Delta$ Rep.\ = (Rep.\ Q2--Q4) $-$ (Rep.\ Q1) em pontos percentuais.
\end{tablenotes}
\end{threeparttable}
\end{table}
